\scanpage{19}
\mspage{15}
% archive sheet number: 18

Du coup, l'intérêt de l'introduction de $\partial\Phi$ \deleted{des} (des sommets-bord de $\Phi$) devient plus problématique, d'autant plus que \deleted{ces sommets} \inserted{l'on n'a} \deleted{\unsure{bord} de $\Phi$ \unsure{ni}} \inserted{plus} évidemment $\partial\Phi=\partial\Phi^*$. \deleted{\unsure{Ci-vant|Ci-contre|à voir}}
\deleted{En fait, on a}
\[
\deleted{\partial\Phi^*\cap\partial\Phi=\partial\Phi\setminus\partial\Phi\cap\partial I}
\]
\deleted{et}
\[
\deleted{\partial\Phi=(\partial\Phi^*\cap\partial\Phi)\amalg\underbrace{\partial\Phi\cap\partial I}_{\substack{\Vert\\\Phi_{\mathrm{int}}\cap\partial I}}}
\]
\[
\deleted{\partial\Phi^*=(\partial\Phi^*\cap\partial\Phi)\amalg\Phi_{\mathrm{is}}\cap\partial I}
\]

On a
\[
\left.\begin{aligned}
\Phi_{\mathrm{is}}&=\Phi^*_{\mathrm{lac}}\\
\Phi_{\mathrm{lac}}&=\Phi^*_{\mathrm{is}}
\end{aligned}\right\}\quad\begin{gathered}\text{\deleted{\unsure{sommets}} lacunaires et s. isolés de $\Phi$ et de $\Phi^*$}\\\text{s'échangent}\end{gathered}
\]
\[
\begin{aligned}
\Phi_{\mathrm{red}}&=\delta\Phi\setminus\delta\Phi\cap\deleted{\unsure{\partial}}\delta\Phi^*\qquad \text{NB. }\Phi^*_{\mathrm{red}}=\partial I\setminus(\partial I\cap\partial\Phi)-\\
&=\bigl(\underbrace{\delta\Phi\setminus\delta\Phi^{*\deleted{*}}}_{\text{sommets évanescents}}\bigr)\amalg\underbrace{(\partial I\cap\partial\Phi)}_{\text{sommets redond}}
\end{aligned}
\]
\[
\deleted{\unsure{\Phi}}\unsure{\partial^*|\delta^*}\Phi=\partial^*\Phi^*\qquad\text{où}\quad\partial^*\Phi=\partial\Phi\setminus\delta\Phi\cap\partial I
\]

Ainsi dans la décomposition canonique
\[
\delta\Phi=\Phi_{\mathrm{red}}\cup\Phi_{\mathrm{is}}\cup\Phi_{\mathrm{lac}}\cup\partial^*\Phi
\]
les $2^{\mathrm{e}}$ et $3^{\mathrm{e}}$ termes s'échangent par dualité, le troisième \unsure{ne varietur|invariant|ne varie pas} par dualité, le premier disparaît par dualité -- mais \deleted{\unsure{réappar}} la partie \deleted{$\partial\Phi\cap\partial I$} $\Phi_{\mathrm{red}}\cap\partial I$ \deleted{\unsure{réapp}} de $\Phi_{\mathrm{red}}$ réapparaît par bidualité: \deleted{\unsure{sommets redon.}} Du coup le terme ``redondant'' pour ces sommets n'est pas tellement adéquat, \deleted{\unsure{on peut les appeler}} si on adopte la terminologie
\[
\text{sommets évanescents}\qquad \Phi_{\text{év}}\qquad\text{old somm.réd. }\Phi_{\mathrm{red}}
\]
on pourrait appeler ces sommets semi-évanescents. On trouve alors une décomposition
